\documentclass[11pt]{article}
\usepackage[margin=0.78in]{geometry}
\usepackage{graphicx}
\usepackage{float}
\usepackage{booktabs}
\usepackage{hyperref}
\usepackage{enumitem}
\setlist{nosep}
\begin{document}
\title{Stage IV-A Mu-Slice Thermodynamic-Boundary Audit}
\author{Report-only audit}
\date{2026-06-28}
\maketitle
\section*{Executive Summary}
The audit uses the frozen Stage IV-A returned dataset \texttt{dataset\_iter025} with 7081 samples. It does not run exact BdG calculations, does not continue active learning, and does not modify thermodynamic or topology labels.

Baseline counts are reproduced: normal=2350, uniform-SC=141, FFLO=4590, trivial=3265, topological=1466, gapless-SC=0, unresolved=0.

The final classification is \textbf{Class A - broad-bin projection artifact}. The recommended next action is \textbf{curve\_extraction\_fix\_only}.

\section*{Direct Answers}
\begin{enumerate}
\item The suspicious wide-bin panels are primarily the high-\(\mu\) bins 0.50\_0.83, 0.83\_1.17, 1.17\_1.50.
\item Broad-bin projection artifact: partial.
\item Narrow fixed-\(\mu\) slices were generated with adaptive half-widths from 0.025 to 0.075.
\item Local normal/SC bracket support was computed with mutual KNN edges. Sampling support limited: False.
\item Thermodynamic margin support is partial; weak-margin fractions are reported in CSV tables.
\item Smooth-fit artifact likelihood is medium.
\item Sampling insufficiency is False.
\item Lower-\(\mu\) edge relation: True; single-band relation: partial.
\item Hard-risk correlation: partial.
\item Overall explanation: Class A - broad-bin projection artifact.
\item Need new exact calculation: False.
\item No exact-tail points are recommended for this specific issue under the current classification. The CSV contains 24 optional traceability candidates if a later decision changes.
\item Stage IV-A remains useful: tFFLO is clearly discovered and the single-band hypothesis is supported, but formal 3D convergence is not yet passed.
\end{enumerate}
\section*{Figures}
\begin{figure}[H]
\centering
\includegraphics[width=0.95\linewidth]{figures/fig01_original_wide_bin_phase_atlas_highlighted.png}
\caption{Original wide-bin phase atlas with the high-mu anomaly bins identified.}
\end{figure}
\begin{figure}[H]
\centering
\includegraphics[width=0.95\linewidth]{figures/fig02_wide_bin_points_colored_by_mu.png}
\caption{Same wide-bin support, colored by actual chemical potential. This tests projection mixing within each wide bin.}
\end{figure}
\begin{figure}[H]
\centering
\includegraphics[width=0.95\linewidth]{figures/fig03_narrow_fixed_mu_exact_label_atlas.png}
\caption{Narrow fixed-mu exact-label atlas without smooth curves.}
\end{figure}
\begin{figure}[H]
\centering
\includegraphics[width=0.95\linewidth]{figures/fig04_narrow_fixed_mu_atlas_with_brackets.png}
\caption{Narrow fixed-mu exact-label atlas with local normal/SC bracket midpoints.}
\end{figure}
\begin{figure}[H]
\centering
\includegraphics[width=0.95\linewidth]{figures/fig05_anomaly_bin_local_bracket_support.png}
\caption{Local normal/SC bracket support in anomaly bins.}
\end{figure}
\begin{figure}[H]
\centering
\includegraphics[width=0.95\linewidth]{figures/fig06_smooth_curve_vs_bracket_supported_curve.png}
\caption{Support-respecting bracket median curves compared with smooth spline diagnostics.}
\end{figure}
\begin{figure}[H]
\centering
\includegraphics[width=0.95\linewidth]{figures/fig07_delta_opt_anomaly_bins.png}
\caption{Thermodynamic order-parameter map in anomaly bins.}
\end{figure}
\begin{figure}[H]
\centering
\includegraphics[width=0.95\linewidth]{figures/fig08_free_energy_margin_anomaly_bins.png}
\caption{Absolute free-energy margin to the normal state in anomaly bins.}
\end{figure}
\begin{figure}[H]
\centering
\includegraphics[width=0.95\linewidth]{figures/fig09_weak_margin_hard_risk_overlay.png}
\caption{Weak-margin and hard-risk overlay near anomalous normal/SC support.}
\end{figure}
\begin{figure}[H]
\centering
\includegraphics[width=0.95\linewidth]{figures/fig10_boundary_reconstruction_sensitivity.png}
\caption{Sensitivity of boundary reconstruction to local KNN/filter choices.}
\end{figure}
\begin{figure}[H]
\centering
\includegraphics[width=0.95\linewidth]{figures/fig11_acquisition_channel_support_by_mu_bin.png}
\caption{Acquisition channel support by mu bin.}
\end{figure}
\begin{figure}[H]
\centering
\includegraphics[width=0.95\linewidth]{figures/fig12_selected_points_anomaly_bins_by_iteration.png}
\caption{Selected points in anomaly bins by iteration and acquisition channel.}
\end{figure}
\begin{figure}[H]
\centering
\includegraphics[width=0.95\linewidth]{figures/fig13_normal_sc_brackets_single_band_projection.png}
\caption{Normal/SC brackets projected over the normal-state single-pair corridor.}
\end{figure}
\begin{figure}[H]
\centering
\includegraphics[width=0.95\linewidth]{figures/fig14_phase_tail_candidates_for_unsupported_arcs.png}
\caption{Report-only phase-tail candidate points for unsupported or weak anomaly arcs.}
\end{figure}
\begin{figure}[H]
\centering
\includegraphics[width=0.95\linewidth]{figures/fig15_final_decision_summary_panel.png}
\caption{Final decision summary panel.}
\end{figure}
\section*{Caveats}
\begin{itemize}
\item Smooth curves and surfaces are report-only diagnostics and are not final physical phase boundaries.
\item The normal-state single-pair corridor is a diagnostic and is not used as a topology label.
\item Candidate points are recommendations only; this script does not submit or run exact calculations.
\item The Stage III hidden-slice reference remains missing from the returned package.
\end{itemize}
\end{document}
